
\section*{Заключение}
\addcontentsline{toc}{section}{Заключение}

В ходе выполнения лабораторной работы:

\begin{enumerate}
    \item Подготовлен кластер из трёх виртуальных машин под управлением
    AlmaLinux 9: узлы \texttt{node1} и \texttt{node2} (10.0.2.11, 10.0.2.12)
    выполняют роль серверов хранения, узел \texttt{mgmt} (10.0.2.10) ---
    роль клиента. Настроено разрешение имён узлов по IP-адресам,
    на серверы установлена серверная часть GlusterFS 11.2, на клиент ---
    клиентская часть \texttt{glusterfs-fuse}.
    \item Узлы \texttt{node1} и \texttt{node2} объединены в доверенный
    пул; оба узла перешли в состояние \texttt{Peer in Cluster (Connected)}.
    Управляющие команды выполнялись на одном узле пула и реплицировались
    на всех пиров.
    \item Созданы и смонтированы на клиенте по сети тома типов
    \texttt{Distributed} и \texttt{Replica}. Установлено: в
    распределённом томе каждый файл хранится целиком на одном брике
    и файлы чередуются между серверами (при отказе узла его файлы
    становятся недоступны); в реплицированном томе каждый файл
    дублируется на обоих серверах, что обеспечивает отказоустойчивость
    ценой двойного расхода дискового пространства.
    \item Изучен механизм квот: на томе \texttt{vrep} включён учёт
    квот, установлен жёсткий лимит 10 МБ на каталог \texttt{/project}.
    После обнуления тайм-аутов применения лимита запись файла 5 МБ
    прошла успешно, а попытка записи следующего файла отклонена
    ошибкой \texttt{Disk quota exceeded}.
    \item Изучен механизм снимков, построенный на тонких томах LVM:
    на серверах развёрнут thin LVM, создан снимок тома, выполнено
    сравнение состояний «до» и «после» (снимок сохранил первую версию
    файла после его изменения в живом томе) и произведено восстановление
    тома из снимка с автоматическим удалением использованного снимка.
\end{enumerate}

Таким образом, все поставленные задачи выполнены в полном объёме,
цель работы достигнута.